% !TEX root = ../main.tex
\chapter{Extending the app}\label{ch:extending}
This chapter condenses the two guides a researcher follows to add to the app,
\code{docs/BLOCK_INTEGRATION.md} and \code{docs/DATA_REGISTRATION.md}, and ends with how to keep
this manual in step with the code. Where this chapter and those guides differ, the guides and the
tests they name are the authority.

\section{Adding a block}\label{sec:extending-block}
\subsection{The contract}
A \term{block}{block} is one \code{AdapterSpec} (\code{Adapters/base.py}) registered from one file,
\code{Adapters/<stage>_<name>.py}, by calling \code{register} (\code{Adapters/registry.py}). The
spec is the whole contract: the executor, the bridge, the pages and the cluster job writer read the
spec and never the function it wraps. Its fields:

\begin{description}
\item[Identity] \code{name} (\code{"<stage>.<short_name>"}, with the stage one of
  \code{Working.recipes.STAGES}: preprocessing, detection, catalogue, comparison, interrogation),
  \code{display_name}, \code{page_name} (what pages call it), \code{category} (its tab in the
  Insert-a-stage window: preprocess, encode, detect, cluster, model or control; an unknown one is an
  error at import), and \code{description}, one paragraph a newcomer can act on.
\item[Types] exactly one \code{input_kind} and one \code{output_kind}, each one of the seven
  interchange types in \code{Working/types/}: \code{signal}, \code{scores}, \code{spanset},
  \code{windowset}, \code{encoding}, \code{grouping}, \code{model}. A block at the start of a chain
  takes \code{signal}; \code{None} is refused at registration.
\item[Settings] a list of \code{ParamSpec}s, each typed (int, float, str or bool), bounded
  (\code{min}, \code{max} or \code{choices}) and with a default. \code{validate_params} fills
  defaults before the recipe is hashed, so typing a default and leaving it out give the same hash
  (\cref{sec:architecture-reproducibility}).
\item[\code{run}] \code{run(x, t, fs, **params, ...)} receives the chain's current signal, its
  absolute time axis and the sampling rate. A block whose input is not a signal also receives the
  previous step's typed output as \code{value}; each declared \term{side input}{side-input} arrives
  under its own name. It returns \code{AdapterResult(output_kind, value=..., meta={...})}:
  \code{value} is the only payload, and anything a page wants besides goes in \code{meta}.
\end{description}

Optional hooks, each read off the spec: \code{side_inputs} (a second typed input: an earlier step,
the root signal or a Library exemplar); \code{estimate} (seconds for this span, or \code{None} when
not calibrated here); \code{max_span_samples} (refuse a longer span before allocating);
\code{recommend} (defaults that depend on the span); \code{derive} (read-only numbers shown beside
the settings); \code{persist} (write the output to disk and register it as an \code{artifacts} row,
so a headless or cluster run produces a browsable result); and \code{known_broken} (registered, but
cannot run here, with the reason). Two keywords on \code{run} itself change how the executor treats
the block: \code{conn} and \code{recording} make it read the live database, so it and every later
step bypass the step cache; \code{units} makes it receive the recording's declared unit, for a block
that prints millivolts (convert with \code{Working.units.read_as_mv}, never a bare factor of 1000;
\cref{sec:architecture-units}).

\begin{watchout}[Index conventions]
\code{x} is the span, so index 0 is the span's first sample. A span set a block emits is
span-relative, and the executor adds the span's start when it writes \code{detections}. A window
set is channel-absolute: its producer adds \code{t[0] * fs}, and a consumer that slices \code{x}
subtracts it. A block that mixes the two draws every window in the wrong place and raises nothing.
\end{watchout}

\subsection{The drawing standard}
A block needs no drawing code of its own. Its picture follows from its types: the \emph{output} type
picks one of seven views, and the \emph{input-to-output} conversion picks a modifier that adds the
evidence (for example, a \code{scores->spanset} block gets its cut drawn on the score curve,
draggable, if it has a float setting named \code{threshold}). The table is
\code{webui/server/views.py}; the client holds one component per view in
\code{webui/client/src/analyse/views/registry.tsx}; every value is serialised through
\code{to_payload} in \code{webui/server/serialize.py}. No view is ever chosen by a block's name. A
block with a conversion no other block has still draws its output type's view, but
\code{tests/test_block_standard.py} fails until a row is added to \code{views.py} and
\code{registry.tsx}. Other conventions a view reads (named \code{meta} keys such as \code{events},
\code{rules}, \code{layers}) are listed in \code{docs/BLOCK_INTEGRATION.md} \S2. The drawing rules
hold for every view: never interpolate, decimate by a minimum--maximum envelope and say so, grey for
no data, and select by sample range, never by window index.

\subsection{Cost}
A block that is not cheap in proportion to the span declares \code{estimate}, returning seconds for
this span or \code{None} when it is not calibrated on this machine, never a guessed constant.
Calibration is per machine: \code{Working/Detection/matrix_profiling/cost.py} and
\code{Working/Preprocessing/window_matrix/cost.py} for the two big ones, and
\code{Working/block_cost.py} for the rest (\code{register_cost_model(name, exponent, run_once)} next
to the spec makes \code{estimate} one line). A block whose cost depends on windows rather than span
length returns \code{None} with that reason, and the chain page then reports the route as unknown.
The estimate is what routes a run to the cluster (\cref{sec:architecture-long-work}). Add a costed
block to \code{COSTED} in \code{tests/test_block_standard.py}.

\subsection{Templates}
A \term{template}{template} is a named, versioned chain, a row of the \code{templates} table. Its
kind follows from the chain's last output type: detection, encoding, training, or interrogation for
a chain ending in an interrogation block. The canonical templates ship as code in
\code{webui/server/templates.py} (\code{CANONICAL}) and are seeded into the table on first start by
\code{seed_canonical}, never overwriting. A multi-stage detector is a template of single-stage
blocks, not one big block, so that each stage can be inspected, swapped and swept: the Dehshibi
detector is \blockref{preprocessing.wavelet_transform} then \blockref{detection.wavelet_summation}
then \blockref{detection.summation_threshold}.

\subsection{The checklist}
\begin{enumerate}
\item Write \code{Adapters/<stage>_<name>.py}: wrap the core function, declare types, settings,
  category, page name and cost, put what a page might want in \code{meta}. No UI imports
  (\cref{sec:architecture-halves}).
\item The file registers itself on import; \code{discover_adapters()} finds it.
\item Test it in \code{tests/test_adapter_<name>.py}: types and category, defaults and bounds, one
  run on a synthetic signal, the error when the typed input is missing. A new output shape needs a
  case in \code{tests/test_webui_serialize.py}; a block in a template needs that template run end to
  end through \code{execute_recipe}.
\item Add its glyph to \code{BY_NAME} in \code{webui/client/src/analyse/glyphs.tsx}.
\item Add a smoke state for its block page to \code{webui/smoke_pages/zz_analyse_run.json}.
\item Run the gate: \code{pytest}, \code{npx tsc -b} and \code{npm run build} in
  \code{webui/client}, and \code{webui/smoke.py} against a running bridge.
\item Document it: run \code{python scripts/manual_tables.py} (the conda Python, from the repository
  root). It adds the block to \code{docs/manual/generated/} and writes stub files
  \code{docs/manual/blocks/<stage>/<short>/card.tex} and \code{entry.tex}. Write the card (plain
  words, for \cref{ch:blocks}) and the entry (exact, for \cref{app:blocks}), following
  \code{docs/manual/AUTHORING.md} and the worked example \blockref{detection.spike_v1}.
\end{enumerate}
Before adding a block, look for one that already does it. Reuse when the meaning matches, not just
the types, and say in the new block's docstring why an existing one would not do.

\section{Adding a kind of data}\label{sec:extending-data}
Nothing on disk is reachable from the app until it is registered. Registration is the same four
verbs for every kind of data (\code{Working/registration/}, routes in
\code{webui/server/registration.py}):

\begin{description}
\item[scan] walks the kind's usual folders and lists every candidate, registered or not, with what
  can be read without loading bulk data (\code{GET /api/registry/{kind}});
\item[check] runs the kind's checks against the file and the database, and never raises: a broken
  file is a failed check that says why (\code{POST /api/registry/{kind}/check});
\item[register] refuses unless every check passes, then writes the row, writes the sidecar manifest
  (\cref{app:data}) and adds an audit entry (\code{POST /api/registry/{kind}/register});
\item[unregister] sets \code{active = 0}; the row and everything that points to it stay
  (\code{DELETE /api/registry/{kind}/{id}}).
\end{description}
Every kind also gets three common checks: the file exists, it is not already registered, and its
hash. Recordings are imported from Settings \textrightarrow{} Datasets (\apppage{settings.datasets}),
models from Settings \textrightarrow{} Models \& registration, and every other kind from a folder's
\ui{scan} on Settings \textrightarrow{} Storage \& backups (\apppage{settings.storage-backups}).

To add a new kind:
\begin{enumerate}
\item Decide its convention (folder, file naming, the manifest fields its producer writes) and add
  it to the kinds table in \code{docs/DATA_REGISTRATION.md}.
\item Add one \code{KindSpec} to the \code{KINDS} dictionary in \code{Working/registration/kinds.py}.
  For a kind stored in \code{registered_artifacts}, the helper \code{_art_kind} fills most fields;
  you write a scan function (facts from headers only) and a check function, which binds the
  recording by content (source file and channel, never a stored row id) with
  \code{_bind_recording}.
\item Test it: scan on a temporary tree, a broken file fails the right check, registering writes the
  right row and sidecar.
\item To offer a \ui{scan} button on Storage \& backups, add a root to \code{get_storage} in
  \code{webui/server/registration.py} and its kind to \code{KIND_OF} in
  \code{webui/client/src/settings/StoragePage.tsx}.
\item Make the producer (a cluster script, a pipeline) write the fields the checks read.
\end{enumerate}

\section{Keeping this manual true}\label{sec:extending-manual}
The manual lives in \code{docs/manual/}, next to the code. Two kinds of content keep it in step.

\begin{description}
\item[Generated] The list of blocks, each block's metadata (names, types, code path, side inputs,
  behaviour flags, the registry's own description) and its settings table come from the adapter
  registry. \code{scripts/manual_tables.py} writes them into \code{docs/manual/generated/}, and
  creates card and entry stubs for any new block; it never overwrites hand-written text. Never edit
  \code{generated/} by hand.
\item[Hand-written] Everything else, including each block's card and entry.
\end{description}

Three commands, all run with the conda Python from the repository root:
\begin{itemize}
\item \code{python scripts/manual_tables.py} after a block is added, removed or changes its
  settings;
\item \code{python scripts/manual_tables.py --check} writes nothing and fails if a generated file is
  stale, a block has no card or entry, or a hand-written folder names a block that no longer exists;
\item \code{python scripts/manual_check.py chapter <file>} or
  \code{... blocks <registry names>}, each with \code{--outdir <folder>}, builds one slice of the book
  and reports LaTeX errors without building the whole thing.
\end{itemize}

When a page is added to the app (\code{webui/client/src/shell/pages.ts}), give it a
\code{\pagesection} with the same id in its workspace's chapter. The rules for writing,
above all ``read the code before you write about it'' and ``never guess'', are in
\code{docs/manual/AUTHORING.md}. A passage not yet written is marked as such, a claim not yet checked
is marked unverified, and the title page counts both.
